\chapter{The Number Splits Into Faces}

\section{Computation names}

A computational name is licensed by a grammar, not by resemblance to a word used elsewhere.  Its record is
\[
 \mathcal N=(\text{name},\text{domain},\text{codomain},\text{route},\text{version},
 \text{inputs},\text{operation},\text{units},\text{budget},\text{stop},
 \text{certificate},\text{projections},\text{erasures},\text{provenance},
 \text{status},\text{forbidden interpretations}).
\]
Two records may share a spelling or a number without sharing an operation.  Conversely, the same operation may admit
several presentations.  The record, rather than resonance, decides what a name owns.

\paragraph{Count.}
A count is \(n\in\mathbb N\) produced by an explicit enumeration or transition rule.  Its units name what is counted:
cells, rungs, partial quotients, candidates, or events.  Failures include invalid domain, exhausted scope, duplicate
membership, overflow, and unresolved enumeration.  Its certificate contains the finite domain or trace, membership
rule, order when relevant, accumulator transitions, scope, and stop.  A count of three partial quotients is not three
grid cells merely because both are written \(3\).

\paragraph{Bit.}
A bit is an element of a declared two-symbol alphabet, such as
\(\mathbb B=\{\mathsf{no},\mathsf{yes}\}\), obtained by a named predicate or encoding map.  Its unit is one symbol;
information in bits requires a separate probability law or coding convention.  Failures include unresolved
predicate, invalid input, unknown encoding, and missing orientation.  Its certificate retains input, predicate,
alphabet, orientation, output, and assumptions.  A Boolean answer alone is neither a probability nor physical spin.

\paragraph{Cost.}
A cost is a nonnegative resource vector accumulated in declared currencies:
\[
 \mathbf c_{i+1}=\mathbf c_i+\mathbf w(s_i\to s_{i+1}),\qquad
 \mathbf c=(\text{requested},\text{spent},\text{remaining}).
\]
The currency may be transitions, arithmetic operations, bit growth, cells visited, checker effort, or elapsed time.
Conversions between currencies require calibration; no scalar total exists without declared conversion weights.
Failures include budget exhaustion, overflow, unknown weight,
and incomparable currencies.  A certificate replays weights and addition along the trace; it does not promote elapsed
time to structural work.

\paragraph{Bracket.}
A bracket is ordered endpoint data \((\ell,u)\) with \(\ell<u\), plus whatever opposed endpoint evidence its claim
requires.  Its units are endpoint units; its width has the same units.  Operations include exact midpoint slicing,
grid conversion, or continued-fraction wall construction, each with a distinct route.  Failures include order,
domain, non-opposed endpoints, arithmetic resource, and unresolved branch.  A certificate owns its route, endpoints,
trace, budgets, and stop, not an unnamed exact point.

\paragraph{Remainder.}
For integers \(a\), \(b>0\), Euclidean division produces the unique pair
\[
                        a=qb+r,\qquad 0\leq r<b.
\]
The remainder is \(r\), in the same integer unit as \(a\) and \(b\).  Division by zero, sign-convention mismatch,
overflow, and an unchecked identity are failures.  Its certificate retains \(a,b,q,r\) and verifies both equality and
range.  A discarded display residue or an empirical leftover may use similar language, but it is not this remainder
without this operation.

\paragraph{Worked example 1: one numeral, different counts (computation).}
The canonical count \(3\) means three whole continued-fraction partial quotients.  That route produces ordered walls
\(17/9\) and \(2\), hence the owned jar \([129.6,137.7]\).  It does not mean three grid cells, dyadic rungs, branch
bits, or any physical count.  Those may also contain the numeral \(3\), but their units, transitions, stops, and
certificates differ.

\paragraph{Worked example 2: a branch bit drives iteration and slicing (computation).}
For ordered bracket \([\ell,u]\), compute \(m=(\ell+u)/2\).  The read bit \(C(m)\) selects either \([m,u]\) or
\([\ell,m]\).  Iteration consumes the bit; slicing discards one half.  The next bracket certificate must retain the
midpoint, bit, retained half, erased half, and endpoint invariant.  Keeping only the bit loses the numerical state;
keeping only endpoints loses the branch history.

\paragraph{Counterexample 3: equal remainders do not identify divisions (computation).}
Both \(17=3\cdot5+2\) and \(23=3\cdot7+2\) have remainder \(2\).  The projection
\(\pi(a,b,q,r)=r\) has many-element fibers.  Equality after \(\pi\) licenses equality of remainders only, not equality
of dividends, divisors, traces, costs, or provenance.

\paragraph{Counterexample 4: equal brackets do not identify processes (computation).}
One trace may receive \([\ell,u]\) as input; another may reach it after several slices.  Equal endpoint pairs certify
the same bracket value under matching units, but not equal trace, route, requested resources, spent resources, or
remaining resources.  Process and cost identity require their own certificate rows.

This manuscript coins \emph{nowtrino} for one scientifically earned computation name:
\[
 \operatorname{nowtrino}(h)=\pi_{\rm now}(h)
   =(\text{current state},\text{current counters},\text{current stop}),
\]
the computation-owned current-summary projection after history slicing.  Its domain is a full finite history \(h\).
Its fiber over a summary contains every history with that same current state, counters, and stop.  The projection
therefore loses prior states, branch order, rejected slices, intermediate costs, and causal route unless separately
retained.  A nowtrino is not a particle, measurement, or physical hypothesis; it is this manuscript's name for a
lossy computational summary.

For example, let \(H_1\ne H_2\) contain different earlier branches but end with the same current state, counters, and
stop.  Then \(\operatorname{nowtrino}(H_1)=\operatorname{nowtrino}(H_2)\).  Equality of summaries is exactly the
projection claim; it cannot recover \(H_1=H_2\).

The spelling \emph{naotrino}, coined from Greek \emph{naos} (temple), is reserved for a later physical apparatus or
model and earns no meaning here.  The etymology records the manuscript coinage; it supplies no physical evidence.  If later
defined, it must state observables, units, calibration, dynamics, uncertainty, and empirical tests.  It must also be
distinguished explicitly from the standard neutrino of particle physics; word resonance establishes no relation.
Any attempted coercion from nowtrino to naotrino or neutrino is rejected as a domain and route mismatch.

Name resolution parses the local declaration and joins on
\((\text{name},\text{domain},\text{route},\text{version})\), requiring a unique row.  The ledger retains origin,
definition, data schema, fiber, and resolution history; a changed meaning requires a new version.  A lexical
concordance finds spellings, not equal semantics or independent evidence.

\section{Physics names}

A physical name is earned by an apparatus, model, observable, units, calibration, uncertainty statement, and data:
\[
\mathcal P=(\text{name},\text{system},\text{apparatus},\text{model},\text{observable},
\text{units},\text{calibration},\text{data},\text{uncertainty},\text{failures},\text{certificate}).
\]
Shared spelling with a computational object supplies none of these.

Computational projections called lower, upper, or binary may be labelled charge, mass, or phase only as an explicitly
declared interpretation.  The labels do not make the projections satisfy \(F=ma\), \(\mathbf F=q\mathbf E\), a phase
model, dimensional laws, or any apparatus response.  Likewise, \(\alpha_c\) is computation-owned: it denotes the
dimensionless ordered-wall jar \([129.6,137.7]\) under its declared route.  It is neither a laboratory determination
nor the fine-structure constant of physics.  Saying that it was generated before it was read records provenance order
only; temporal order does not validate a physical interpretation.

\paragraph{Mass.} In an inertial experiment, \(F=ma\).  Calibrated force and acceleration observations estimate
\(m\) in kilograms.  Alignment, friction, sensor bias, timing, model inadequacy, and nonidentifiability are failures.
The certificate retains calibration, observations, fit, residuals, uncertainty, and validity range.

\paragraph{Charge.} Charge \(q\), in coulombs, may be inferred from \(\mathbf F=q\mathbf E\).  Apparatus must measure
force and field orientation.  Background fields, leakage, sign ambiguity, saturation, and nonuniformity limit the
claim.  Geometry, field map, zeroing, sign convention, data, corrections, and uncertainty belong in the certificate.
The equation alone may be nonidentifying: for any admissible \(\lambda\ne0\), the pairs
\[
                         (q,\mathbf E),\qquad(\lambda q,\mathbf E/\lambda)
\]
produce the same \(\mathbf F\).  The report fiber over one force therefore contains multiple charge--field pairs
unless another calibrated channel fixes the scale.  Dimensions are necessary consistency checks, not sufficient
identification evidence.

\paragraph{Phase.} For \(y(t)=A\cos(\omega t+\phi)\), phase \(\phi\) is in radians modulo \(2\pi\), relative to a
declared clock and oscillator.  Drift, aliasing, low amplitude, cycle slips, and reference changes are failures.
Circular uncertainty and reference epoch are required.

\paragraph{Field.} A field such as \(\mathbf E:\Omega\times\mathbb R\to\mathbb R^3\), in volts per metre, is sampled
only at finite locations, times, orientations, and bandwidths.  Interpolation is a declared model.  Drift, spatial
averaging, boundaries, missing samples, and residuals limit the certificate.

\paragraph{Coupling.} In \(y_i=gx_i+\epsilon_i\), coupling \(g\) has units \(y/x\), unless normalization makes it
dimensionless.  Both channels, excitation range, error covariance, backgrounds, and model selection require
calibration.  Raw pairs, estimator, residuals, uncertainty, and domain are retained.

\paragraph{Worked measurement 1: mass (physics).} Apply three calibrated forces to a carriage and measure repeated
accelerations.  A weighted linear fit estimates \(m\), and residuals test the model.  Numerical iteration may solve
the fit and slicing may form an interval, but convergence alone is not a mass certificate: friction, timing, force
calibration, and the inertial model remain physical requirements.

\paragraph{Worked measurement 2: phase (physics).} Sample a sinusoid against a traceable clock.  Numerical search may
iterate over \(\phi\) and retain a likelihood slice.  Changing the reference epoch shifts phase without changing the
waveform, so the physical name includes its reference and modulo convention.

\paragraph{Counterexample 3: a slope is not yet a coupling (physics).} Two digitized channels yield slope \(2.0\).
If one channel used millivolts rather than volts, the physical slope changes by \(10^3\).  The numeral cannot name a
coupling until units, calibrations, model, and uncertainty are supplied.

\paragraph{Counterexample 4: samples are not a continuous field (physics).} Equal readings at two probes do not
exclude variation between them.  Interpolation or smoothness can fill the gap only as a model assumption.  Numerical
slicing bounds parameters within that model; it cannot observe unsampled points.

The manuscript spelling \emph{naotrino} remains reserved.  No concrete apparatus, model, observable, units,
calibration chain, uncertainty analysis, or discriminating dataset is supplied here.  It is not the computation-owned
nowtrino and not the standard neutrino of particle physics.  A later coinage requires a complete \(\mathcal P\) record
and evidence against alternatives.

Numerical acquisition preregisters the observable, sampling plan, blinding rule, preprocessing, estimator, protocol
and version, and budgets before it parses calibrated samples.  It validates units, iterates the declared estimator,
slices a parameter region, and reports pass, fail, or unresolved checks.  Raw and rejected samples are preserved with
their rejection reasons; blinding is lifted only by the preregistered rule.  Stops include acquisition budget, missing
calibration, unit mismatch, saturation, domain failure, nonidentifiability, iteration budget, overflow, tolerance,
complete, and model rejection.  A numerical stop is not a physical failure absent an explicit apparatus relation.

Different named estimates can share data, calibration standards, preprocessing, environmental disturbances, or fitted
parameters.  Their covariance and cross-covariance must therefore be estimated or bounded under the joint acquisition
model.  Distinct names, equations, or report columns establish no statistical independence.  Treating correlated
channels as independent can shrink uncertainty without new evidence.

The physical report must retain apparatus, observable, units, calibration, sampling, model, residuals, uncertainty,
covariance, shared preprocessing, and validity range.  A changed calibration, model, or apparatus starts a new report;
none of these records follows from a computational name.

\section{Same number or same process}

Comparison begins by declaring its relation.  For reduced rationals,
\[
a/b =_{\rm exact} c/d\Longleftrightarrow ad=bc,
\qquad
x =_\theta y\Longleftrightarrow q_\theta(x)=q_\theta(y)
\]
defines exact equality and equality under declared quantizer \(q_\theta\).  Exact equality implies resolution equality
only inside its domain; resolution equality does not imply exact equality.

A process report is
\[
P=(\text{inputs},\text{algorithm},\text{transitions},\text{trace},\text{cost vector},
\text{budgets},\text{stop},\text{configuration},\text{provenance}).
\]
Process equivalence requires comparison maps for every component: compatible inputs, transition semantics, trace
correspondence, cost currencies, budgets, stops, configuration, and provenance.  Output equality is only one row.
Shared names imply none of them.

Value and process provenance remain separate:
\[
\text{raw}_i\xrightarrow{q_{\theta_i}}\text{label}_i,
\qquad
\text{input}_i\xrightarrow{P_i}\text{output}_i.
\]
A value comparison map cannot fill absent input, transition, trace, or cost maps.

\paragraph{Worked comparison 1: direct and stepped endpoints (computation).}
At scale \(10^{18}\),
\[
d=137011290548979455469,\qquad s=137011290753751157155.
\]
The integers differ by \(|d-s|=204771701686\), while a declared coarse display maps both to \(137.011\).  Exact
inequality, declared-resolution equivalence, and process difference are separate rows.  The shared formatted string is
a formatter collision, not equality.  Direct and stepped routes have different transitions and costs.

A declared inverse-report map takes continued-fraction walls \(17/9\) and \(2\) to jar \([129.6,137.7]\).  The map and
calibration belong in the certificate.  Wall rationals and report endpoints inhabit different numerical domains; they
are not the same numbers.

\paragraph{Worked comparison 2: same operation, different inputs (computation).}
Division of \(17\) by \(5\) and \(23\) by \(7\) uses the same operation and returns remainder \(2\).  Inputs and quotient
identities differ.  The runs share an algorithm schema but are not the same process.

\paragraph{Counterexample 3: same value, different algorithms (computation).}
One route reduces \(18/12\) by remainders; another reads \(3/2\) from a table.  Exact outputs agree, but traces,
failures, costs, and certificates differ.  Value equality passes while process equivalence fails or is unresolved.

\paragraph{Counterexample 4: shared names provide no bridge.}
A nowtrino is a computation-owned summary.  A naotrino is a reserved possible physical name, while a standard neutrino
belongs to established physics.  No comparison map, calibration, or model bridge is supplied.  Resemblance cannot
construct one.

Comparison joins relation, domain, route, version, units, and configuration before testing value, resolution, trace,
cost, and provenance.  Endpoint projection erases trace and cost; display projection erases residue; current summary
erases history; and an unkeyed name erases domain and route.

Same-name means resolution to one governed ledger row---name, domain, route, version, and status---not equal spelling.
Typed value paths are
\[
\mathrm{Input}_i\to\mathrm{Report}_i\xrightarrow{\kappa_i}C\xrightarrow{q_\theta}\mathrm{Label}_i.
\]
Each \(\kappa_i\) converts to common comparison domain \(C\).  Process comparison instead uses an explicit trace/state
relation, not \(\kappa_i\).  Byte-for-byte trace equality is strong; a weaker relation must state which states,
transitions, stuttering steps, and costs it identifies.  Without one, process equivalence is unresolved.

\paragraph{Worked comparison 5: finite three-state closure (computation).}
Let \(S=\{\mathsf{minus},\mathsf{zero},\mathsf{plus}\}\) and declare the deterministic rotation
\[
\mathsf{minus}\mapsto\mathsf{zero}\mapsto\mathsf{plus}\mapsto\mathsf{minus}.
\]
Finite case reduction certifies for every \(s\in S\) that three rotations return \(s\), while one and two rotations do
not.  The displayed finite table suffices under the declared transition.  Its scope is exactly these three states.
Closure does not make an unrelated residue a one-step
invariant or fixed point.  Only a separately defined residue map with a proved preservation law could do that.
After the three-step return, any read defined solely from the returned state agrees with its initial state-derived
read.

\paragraph{Worked comparison 6: reproducible tuples are not cycle causation (computation).}
Each of two governed runs returns the exact tuple
\[
                         (573,573,552),
\]
with residual \(552-573=-21\) and scaled inverse
\[
                         137011290548979455469.
\]
Tuple equality, residual arithmetic, and scaled-inverse equality are separate rows.  The receipt records checker and
compiler releases, snapshot hash, configuration, deterministic inputs, and declared nondeterministic inputs for both
runs.  Equality is \emph{read}; configuration sameness is posited and validated against those records.  Passing rows
establish reproducibility under recorded assumptions, not that the three-state cycle caused the tuple or residual.

Charge labels, RGB labels, or \(a,b,c\) labels exhibit the same three-fold organization only when declared maps to
\(S\) are supplied and their transition-preservation squares are checked.  Without those maps, the resemblance is an
interpretation, not process or value equality.

\paragraph{Descent experiment 7: sign projection versus signed residue (finite apparatus audit).}
Use the live finite apparatus surface only as an executable computation.  Its input domain is the product
\[
 \mathcal X=\mathcal P\times\mathcal V,
\]
where \(\mathcal P\) contains the declared four-node cubic paths and \(\mathcal V\) contains declared two-leg
variations.  The first transform is the mixed second-difference read
\[
 h:\mathcal X\to\mathbb Z,\qquad h(p,v)=\operatorname{holonomy}(p,v),
\]
and the second is the sign classifier
\[
 \sigma:\mathbb Z\to
 \{\mathsf{matter},\mathsf{antimatter},\mathsf{doesntMatter}\}.
\]
For the checked loop variation \(v_{12}\), the finite table reads
\(h(p_{\rm flat},v_{12})=-1\) and \(h(p_{\rm tilted},v_{12})=+1\).  For the held-right-leg variation \(v_{21}\) on
the tilted path it reads \(h(p_{\rm tilted},v_{21})=0\).  Applying \(\sigma\) yields the actual source constructor
\texttt{Charge.doesntMatter} at zero; \emph{neutral} is only a display gloss.  This is
the exact pointwise computation of the source definitions on these named inputs: path and variation to signed
integer, then integer to label.  It is
not a transform from a physical particle to a measured charge.

Define the direct event read \(E(p,v)\) by evaluating the program's event constructor.  The comparison test is the
commuting triangle
\[
                       E=\sigma\circ h
\]
on the three enumerated inputs.  Pointwise equality passes for the returned labels.  The converse square fails:
there is no inverse \(\sigma^{-1}\), because all positive integers share the antimatter label and all negative
integers share the matter label.  The residual of sign projection is the erased magnitude and path/variation
provenance.  Equal labels therefore support neither equal residues nor equal processes.

This is a pointwise specialization of the live definitions on three named inputs, not an exported run, trace schema,
certificate, or stop vocabulary.  Its distinct ownership lesson is that integer equality, constructor equality, and
path--variation equality are separate rows.  Chapter 9 carries the canonical transform and adversarial audit.

The integer evaluations are live-definition results; the triangle is a manuscript restatement, not an exported
report.  Its physical names add no calibrated observation.  Ownership remains with these three input rows and the
separate equality relations above.
